\begin{formulario}{T-6}
  \begin{formularioTSeis}

    \proyectoTSeis{
      nombre       = Sistema Híbrido de Telemetría\, Trazabilidad y Control de Flota,
      cliente      = Transportes del Sur Ltda.,
      industria    = Transporte terrestre interurbano de carga pesada y distribución troncal,
      periodo      = 2023--2024 (14 meses de implementación; en operación continuada),
      monto        = Rango histórico superior a 50.000 Unidades de Fomento,
      alcance      = Provisión e instalación de pasarelas embarcadas CANbus J1939; desarrollo de plataforma cloud de despacho\, control de geocercas dinámicas y monitoreo en tiempo real.,
      arquitectura = Híbrida: Edge computing en cabina + nube pública Microsoft Azure,
      servicio     = 99{,}5\,\% mensual en disponibilidad de plataforma 24/7/365,
      volumen      = 340 tractocamiones activos; $\approx 88.000$ viajes anuales; $> 6{,}8$ millones de eventos/día,
      rol          = Contratista Principal (100\,\% ingeniería\, desarrollo y soporte),
      contacto     = Marcelo Iturra Valenzuela\, Gerente de Operaciones
    }

    \proyectoTSeis{
      nombre       = Plataforma IoT y Telemetría Resiliente con Arquitectura Offline-First,
      cliente      = AgroFrutícola Los Andes S.A.,
      industria    = Agroindustria\, faenas remotas y logística de exportación en frío,
      periodo      = 2022--2023 (12 meses de implementación; en operación continuada),
      monto        = Rango histórico entre 30.000 y 45.000 Unidades de Fomento,
      alcance      = Despliegue de concentradores telemáticos locales on-premise\, integración de 1.500 sensores de frío PT100 y sincronización bidireccional offline-first.,
      arquitectura = Híbrida: Concentradores locales on-premise en plantas + nube Microsoft Azure,
      servicio     = 99{,}2\,\% mensual; tolerancia a operación desconectada de 72 horas,
      volumen      = 310 unidades de transporte y frío; 12 faenas remotas; $\approx 82.000$ despachos/año,
      rol          = Contratista Principal (100\,\% ingeniería\, integración y despliegue),
      contacto     = Paula Concha Morales\, Jefa de Tecnologías de Información
    }

    \proyectoTSeis{
      nombre       = Plataforma Analítica Centralizada y Gestión Logística en Alta Disponibilidad,
      cliente      = Logística y Distribución Multimodal Bicentenario S.A.,
      industria    = Logística portuaria\, distribución multimodal y transferencia de carga,
      periodo      = 2024--2025 (15 meses de implementación; en operación continuada),
      monto        = Rango histórico superior a 60.000 Unidades de Fomento,
      alcance      = Arquitectura cloud escalable en Azure con microservicios e ingesta de eventos de alta velocidad; reportería analítica y despacho automatizado.,
      arquitectura = Híbrida: Nube Microsoft Azure articulada con servidores de borde en terminales,
      servicio     = 99{,}6\,\% mensual garantizado contractualmente con monitoreo continuo,
      volumen      = 390 tractocamiones y portacontenedores; $\approx 98.000$ viajes anuales; $> 12$ millones de eventos/día,
      rol          = Contratista Principal (100\,\% arquitectura\, software e implantación),
      contacto     = Rodrigo Baeza San Martín\, Subgerente Corporativo de Sistemas
    }

  \end{formularioTSeis}
\end{formulario}
